%% LyX 2.5.1 created this file.  For more info, see https://www.lyx.org/.
%% Do not edit unless you really know what you are doing.
\documentclass[12pt,english]{article}
\usepackage{mathpazo}
\usepackage[T1]{fontenc}
\usepackage[latin9]{inputenc}
\usepackage{float}
\usepackage{geometry}
\geometry{verbose,tmargin=1cm,bmargin=2.7cm,lmargin=2.7cm,rmargin=2cm}
\pagestyle{empty}
\usepackage{setspace}
\usepackage[authoryear]{natbib}
\onehalfspacing

\makeatletter
\@ifundefined{date}{}{\date{}}
%%%%%%%%%%%%%%%%%%%%%%%%%%%%%% User specified LaTeX commands.
% general
\usepackage[titletoc]{appendix}
\usepackage{graphicx}
\usepackage{placeins}
\usepackage{tikz}

% algorithm
\usepackage[ruled,vlined,linesnumbered]{algorithm2e}
\IncMargin{2.5cm}
\DecMargin{2cm}
\usepackage{fullwidth}
\usepackage{enumitem}
\setlist{leftmargin=1.7cm}

% tables
\usepackage{tabularx, siunitx, multirow, booktabs}
\begingroup
% Allow `_` and `:` in macro names (LaTeX3 style)
\catcode`\_=11
\catcode`\:=11
% Internal code of `S`
\gdef\tabularxcolumn#1{%
    >{\__siunitx_table_collect_begin:Nn S{} }%
    p{#1}%  <- this is different (is `c` in normal `S`)
    <{\__siunitx_table_print:}%
}
\endgroup

% figures
\usepackage{subfig}
\usepackage{caption}
\captionsetup[subfloat]{position=top}

% footnotes
\setlength{\skip\footins}{1cm}
\usepackage[hang,splitrule]{footmisc}
\setlength{\footnotemargin}{0.3cm} %.5
\setlength{\footnotesep}{0.4cm}

% code
\usepackage{xcolor}
\usepackage{listings}

\definecolor{codegray}{rgb}{0.5,0.5,0.5}
\definecolor{background}{HTML}{F5F5F5}
\definecolor{keyword}{HTML}{4B69C6}
\definecolor{string}{HTML}{448C27}
\definecolor{comment}{HTML}{448C27}

\usepackage{inconsolata}
\lstdefinestyle{mystyle}{
    commentstyle=\color{comment},
    keywordstyle=\color{keyword},
    stringstyle=\color{string},
    basicstyle=\ttfamily,
    breakatwhitespace=false,         
    breaklines=true,                 
    captionpos=b,                    
    keepspaces=true,                                    
    numbersep=5pt,                  
    showspaces=false,                
    showstringspaces=false,
    showtabs=false,
    tabsize=4,
	showlines=true
}

\lstset{style=mystyle}

% manual
\usepackage{enumitem}
\setlist[enumerate]{leftmargin=1cm}
\setlist[itemize]{leftmargin=0.5cm}
% Added by lyx2lyx
\setlength{\parskip}{\smallskipamount}
\setlength{\parindent}{0pt}

\makeatother

\usepackage{babel}
\begin{document}
\title{\textsc{Assignment III: Your Model}}

\maketitle
\vspace{-3mm}\thispagestyle{empty}\textbf{Vision: }This project teaches
you to construct a heterogeneous agent model and use it to analyze
an economic question.
\begin{itemize}
\item \textbf{Tasks:} Your answer must contain
\begin{enumerate}
\item A self-contained verbal and mathematical description of a heterogeneous
agent model of your choice.
\item A detailed step-by-step description of how you would use the model
to answer an economic question of your choice.
\end{enumerate}
The model can be an extension of a model you have seen in the course
(e.g. in an assignment) or elsewhere.

The good answer provides a clear exposition of the model and the proposed
analysis as well as a creative application of the tools learned in
the course.
\item \textbf{Hand-in: }Upload a single pdf-file on Absalon (and nothing
else). Maximum of 5 pages.
\item \textbf{Deadline (hand-in):} 10th of December 2025
\item \textbf{Deadline (peer feedback):} 14th of December 2025
\item \textbf{Exam: }Your Assignment III will be a part of your exam portfolio.\\
You can incorporate feedback before handing in the final version.\\
\emph{Note: Peer-feedback is a mandatory part of the assignment.}
\end{itemize}

\end{document}
